\begin{claim}{C.49}{Proved}{Lie contractions on closures with three boundary foliations}
Let $f_t$ be an ambient-null based loop in
$\operatorname{Homeo}^+(\R^2)$, let $G$ be its generated group,
and let $H$ commute with every $f_t$.
For every open $G$-orbit $O$ having at least three distinct actual
lifted boundary-quotient fiber partitions of C.42, the group
generated by the based lifts of $f_t|_O$ to the universal covering
$X\to O$ is closed in $\operatorname{Homeo}(X)$. With this actual
topology it is one of
\[
 \R^2,\qquad \operatorname{Aff}^+(\R),\qquad
 \mathcal A_2=\R^2\rtimes\R ,
                                                        \tag{C.49.1}
\]
where the last factor acts on $\R^2$ by positive scalar
dilations. Each group is simply connected and has a globally
invertible exponential map.

There is a jointly continuous based contraction on the whole
$\overline O$, with jointly continuous inverse, obtained from
these actual lifted groups by fractional powers.
These contractions preserve every actual boundary-quotient
equation, agree on intersections of such orbit closures and
with the contraction of C.37 on $X_2$, and are transported to
one another by $H$.
They consequently give a continuous homeomorphism family on
every finite union of these closures and $X_2$.
The conclusion does not assert continuity across an arbitrary
infinite accumulating family of the new three-dimensional
cases.
\end{claim}

\begin{proof}
\emph{The actual based source lifts are closed loops and act transitively.}
Choose three partitions and lifts $F_i:X\to\R$ as in C.47.
Lift the isotopy $f_t|_O$ from the identity to
$\widetilde f_t\in\operatorname{Homeo}(X)$.
The covering homotopy property gives continuity for maps and
inverses in their compact-open topology.
Every $\widetilde f_t$ commutes with every deck transformation:
the commutator is a continuous deck-valued function of $t$
and is the identity at $t=0$.
Its target lifts satisfy
\[
                    F_i\widetilde f_t=b_{i,t}F_i .
                                                        \tag{C.49.2}
\]
For a real target, $b_{i,1}=\id$ directly. For a circle target,
the exact two-point winding comparison of C.29 makes the
based real lift return to the identity, since the original
plane loop is ambient-null. Thus $b_{i,1}=\id$ for all $i$.
The injectivity of $(F_1,F_2)$ from C.45 now gives
$\widetilde f_1=\id$.

Let $N_{\mathrm{alg}}=\langle\widetilde f_t:0\leq t\leq1\rangle$.
This actual subgroup is path connected: each finite word is
joined to the identity by simultaneously shortening its
factor paths. Put
\[
 P=\{\widetilde f_t:0\leq t\leq1\}
       \cup\{\widetilde f_t^{-1}:0\leq t\leq1\}.
 \]
It is compact and contains the identity.
The source deck group $\Delta$ is countable, because a fiber
of a covering of a second-countable surface is a discrete
countable subset of its second-countable covering plane.
For $x\in X$, transitivity of $G$ on $O$ gives
\[
                       X=\bigcup_{\delta\in\Delta}
                            \ \bigcup_{n\geq1}\delta P^n x .
                                                        \tag{C.49.3}
\]
All sets in this countable union are compact.
Baire's theorem gives interior to one of them, hence to
$N_{\mathrm{alg}}x$. Transitivity on that orbit then makes
the entire orbit open.
Every other orbit is a deck translate of it, since the
generated lifted group commutes with $\Delta$.
These disjoint open orbits partition the connected plane;
there is only one. Thus $N_{\mathrm{alg}}$ is transitive
on $X$, not just on its projection $O$.

\emph{The connected three-web groups and their exact actions.}
Let $\mathcal W$ be the full group in C.47 and write
$W=\mathcal W^0$. The group $W$ is open in $\mathcal W$.
The open orbit map makes every $W$-orbit open in $X$;
normality of $W$ and connectedness of $X$ imply transitivity.
For a point $x$, let $J=W_x$.
The Lie principal bundle $W\to W/J=X=\R^2$ shows that
$J$ is connected: the exact sequence on components has
$\pi_1(X)=0$ and $\pi_0(W)=0$.

If $\dim W=2$, then $J$ is discrete and connected, hence
trivial. Evaluation identifies $W$ homeomorphically with
the simply connected plane. Its Lie algebra and group are
the two explicit possibilities $\R^2$ and
$\operatorname{Aff}^+(\R)$ in C.36.
Since $N_{\mathrm{alg}}\subset W$ is already transitive
and this action of $W$ is simply transitive, one has
$N_{\mathrm{alg}}=W$.

Suppose $\dim W=3$. Here C.47 puts us in the case where all
three pair images are full planes.
The connected one-dimensional isotropy $J$ is a subgroup
of the ordered-loop automorphism group in that proof.
It cannot be a circle. Indeed, every compact group of
increasing line homeomorphisms is trivial: the compact orbit
of any point has a largest element, which all group elements
must fix, so transitivity on that orbit makes it a singleton.
Thus $J\cong\R$. The principal bundle over the contractible
plane now gives $\pi_1(W)=0$.

We identify its Lie algebra and homogeneous action explicitly.
Let $\mathfrak w=\operatorname{Lie}(W)$ and
$\mathfrak j=\operatorname{Lie}(J)$.
For each $i$, the stabilizer $K_i$ of $F_i(x)$ in the
induced real target action is a closed subgroup containing
$J$. The open orbit map identifies $W/K_i$ with the actual
real target. The same principal-bundle component argument
makes $K_i$ connected, and its Lie algebra $\mathfrak k_i$
has dimension two.
The three $\mathfrak k_i$ are distinct: otherwise the
corresponding connected subgroups, and hence the quotient
fibers through $x$, would coincide, contradicting C.45.

For $0\ne Z\in\mathfrak j$, the induced linear map
$\operatorname{ad}Z$ on $\mathfrak w/\mathfrak j$ preserves
the three distinct lines $\mathfrak k_i/\mathfrak j$.
A real linear map on a two-dimensional space preserving
three distinct lines is scalar; write that scalar $\lambda$.
If $\lambda=0$, then
$[\mathfrak j,\mathfrak w]\subset\mathfrak j$, so
$\mathfrak j$ is an ideal.
Its connected subgroup $J$ is then normal in $W$.
As a normal point stabilizer in a transitive action it fixes
all of $X$, contrary to the faithful inclusion
$W\subset\operatorname{Homeo}(X)$.
Therefore $\lambda\ne0$.

The adjoint map on $\mathfrak w$ has the eigenspace
decomposition
\[
 \mathfrak w=\mathfrak j\oplus V,\qquad
 \dim V=2,\qquad [Z,v]=\lambda v\quad(v\in V).
 \]
For $v,w\in V$, the derivation identity gives
$[Z,[v,w]]=2\lambda[v,w]$.
The value $2\lambda$ is absent from the spectrum
$\{0,\lambda\}$, so $[V,V]=0$.
Replacing $Z$ by $Z/\lambda$ identifies the Lie algebra
with scalar homotheties:
\[
              [Z,v]=v,\qquad [V,V]=0.                  \tag{C.49.4}
\]
Since $W$ is simply connected, the corresponding group
is exactly $\mathcal A_2$, with multiplication
\[
        (a,v)(b,w)=(a+b,\ v+e^a w),
        \qquad a,b\in\R,\quad v,w\in\R^2.               \tag{C.49.5}
\]
Under this isomorphism $J$ is the dilation subgroup
$\{(a,0)\}$. The orbit identification
$W/J\to X$ therefore gives an actual equivariant
homeomorphism of the source plane with the standard action
\[
                         z\longmapsto v+e^a z .
                                                        \tag{C.49.6}
\]

\emph{The generated subgroup is closed, rather than merely analytic.}
Yamabe's theorem makes the path-connected subgroup
$N_{\mathrm{alg}}\subset W$ an analytic Lie subgroup
\cite[pp.~13--14]{C_Yamabe1950}.
This statement initially gives its analytic subgroup
structure; it does not identify that structure with the
induced topology. We now establish the latter point by
classifying the transitive subgroups of (C.49.6).

A transitive analytic subgroup has dimension at least two.
In dimension two its Lie algebra is either the translation
plane $V$, or has a nonzero projection to the dilation
coordinate. In the latter case it is
\[
                    \R(Z+v_0)\oplus L,                 \tag{C.49.7}
\]
where $L\subset V$ is a one-dimensional translation subspace.
On the quotient line $V/L$, the generator $Z+v_0$ has flow
\[
                 \bar z\longmapsto
                       e^t\bar z+(e^t-1)\bar v_0 .
 \]
It fixes $-\bar v_0$, while $L$ acts trivially there.
Thus the analytic subgroup in (C.49.7) preserves the proper
affine line $-v_0+L$ in the source plane, contradicting
transitivity. The remaining two-dimensional subgroup is
exactly all translations $V$.
A three-dimensional analytic subgroup is open and equals
the connected group $\mathcal A_2$.
Hence, in this case,
$N_{\mathrm{alg}}=V$ or $N_{\mathrm{alg}}=\mathcal A_2$.

Both are closed with their actual homeomorphism-group
topology. One evaluation recovers a translation vector.
For a homothety, evaluation at two fixed distinct points
recovers its positive scale and translation continuously;
a limit that is still a homeomorphism has positive scale.
Together with the dimension-two case already treated,
this proves that $N=N_{\mathrm{alg}}$ is closed and is
one of (C.49.1), with its induced topology.
It also proves that choosing a different triple of
partitions does not change $N$: it is the same subgroup
generated by the same based lifted loop.

For the semidirect products in (C.49.1), with vector
dimension one or two, the exponential and its inverse
are given by
\[
 \exp(u,v)=
   \left(u,\frac{e^u-1}{u}\,v\right),
 \qquad \frac{e^u-1}{u}=1\quad(u=0).                    \tag{C.49.8}
\]
The displayed scalar factor is positive everywhere.
For $\R^2$ the exponential is the identity.
This proves the asserted global exponential coordinates
without an unspecified choice of logarithm.

\emph{The exact Polish quotient extending the action to the closure.}
Put $L=\overline G$ in the plane homeomorphism group.
It preserves $O$ by the forward-and-inverse boundary argument
of C.36. Define the closed Polish subgroup
\[
 E_O=\{(l,n)\in L\times N:
                  p n=(l|_O)p\}.                       \tag{C.49.9}
\]
Projection $E_O\to N$ is onto: each $n\in N$ is an actual
finite word in the based lifted loop, so the identical
word in the original plane loop supplies its $l\in G$.
This continuous surjective homomorphism of Polish groups
is open.
Its kernel consists of pairs $(l,1)$ with $l$ fixing $O$
pointwise, and therefore fixing $\overline O$ pointwise.
Consequently
\[
 N\times\overline O\longrightarrow\overline O,\qquad
     (n,z)\longmapsto lz\quad\bigl((l,n)\in E_O\bigr)     \tag{C.49.10}
\]
is well defined and jointly continuous.
Indeed, the product of the open projection with
$\id_{\overline O}$ is an open quotient, and the continuous
$E_O$-action factors through it.
This argument keeps the exact source covering and the
possibly nontrivial global restriction kernel.

Every $n$ has its chosen preimage in $G$, so (C.49.10)
preserves every original $G$-orbit in $\overline O$ and
has on that orbit exactly the image of $G$.
It commutes with the original setwise $H$-stabilizer of
$O$, since those preimages commute with that stabilizer.
All source deck transformations commute with $N$ because
they commute with its generating based lifts, so the
restriction of (C.49.10) to $O$ is also the exact descended
source action.

The required closure homotopy is now
\[
 R^O_{s,t}
   =\exp\bigl((1-s)\log\widetilde f_t\bigr)
           \quad\hbox{acting by (C.49.10)}.              \tag{C.49.11}
\]
The source loop is closed and based, as proved above.
Thus (C.49.11) starts with $f_t|_{\overline O}$, ends
with the identity, and fixes both based edges.
The action and group inversion give joint continuity
of the family and its inverse.

\emph{All actual boundary quotients and common boundary images.}
Consider any nontrivial nonopen $G$-orbit
$B\subset\overline O$, with its intrinsic line or circle
topology from C.28. The action (C.49.10) is transitive
on its underlying set. Its point stabilizer $N_z$ is closed.
The continuous bijective orbit map $N/N_z\to B$ has
one-dimensional domain.
Dimension zero is excluded by connectedness and nontriviality.
If its dimension were at least two, restricting to a small
two-dimensional coordinate disk and applying invariance
of domain in the ambient plane would give interior to
$B$, making this whole transitive orbit open.

The bijection agrees with the intrinsic topology.
A compact interval in the homogeneous one-manifold maps
to an embedded arc. On a non-circular orbit that arc is
the unique connecting arc of C.28; on a circle the same
statement holds in a sufficiently small interval.
Its interior is therefore an intrinsic interval
neighborhood. This proves that the orbit identification
is a homeomorphism, and that the $N$-action is continuous
in the intrinsic topology as well.

Every group in (C.49.1) is metabelian.
Its image on $B$ is path connected and equals the actual
$G|_B$. The exact metabelian image argument of C.37
therefore identifies this common image, with its induced
intrinsic compact-open topology, as
\[
 M_B=
 \begin{cases}
   \R\text{ acting by translations, or }
        \operatorname{Aff}^+(\R),&B\text{ a line},\\
   \R/\Z\text{ acting by rotations},&B\text{ a circle}.
 \end{cases}                                           \tag{C.49.12}
\]
These groups are closed in the intrinsic homeomorphism
group: one evaluation recovers the translation or
rotation, and two recover the affine parameters.
The action supplies a continuous surjective Lie
homomorphism
\[
                            \rho_B:N\to M_B.            \tag{C.49.13}
\]
In particular every actual boundary quotient retains its
original equivariance equation under the closure action.
To verify this directly for
$\Phi_z(lx)=lz$, take the preimage $g\in G$ of $n\in N$.
The equality $\Phi_z(gy)=g\Phi_z(y)$ follows from the
defining quotient formula; replacing $g$ by another
preimage changes neither side because its restriction
kernel fixes $\overline O$. Thus
\[
            \Phi_z(R^O_{s,t}y)
                   = R^B_{s,t}(\Phi_z(y))               \tag{C.49.14}
\]
with the same original line or circle target and its
covering equation.

Since $N$ is simply connected, (C.49.13) has a unique
based lift $\widetilde\rho_B:N\to\widetilde M_B$.
The latter group is $\R$ or $\operatorname{Aff}^+(\R)$,
whose exponential is globally invertible.
The path $\widetilde\rho_B(\widetilde f_t)$ is precisely
the based lift $\widetilde b_t$ of the original
$f_t|_B$. Naturality and injectivity of the exponential
give the exact identity
\[
 \widetilde\rho_B
       \left(\exp\bigl((1-s)\log\widetilde f_t\bigr)\right)
       =\exp\bigl((1-s)\log\widetilde b_t\bigr).          \tag{C.49.15}
\]
This defines $R^B_{s,t}$ in (C.49.14) after projection.
On a point orbit the action is the identity.

\emph{Intersections, operators, and the finite-union conclusion.}
For two orbit closures, their intersection consists of
nonopen $G$-orbits. On each nontrivial one, both images
(C.49.12) are the identical group $G|_B$ with the identical
intrinsic topology. The two based target-cover paths
are the same lift of $f_t|_B$. Formula (C.49.15)
therefore makes the closure contractions equal there.
It is also exactly the boundary formula in C.36--C.37,
so it agrees on every nonopen orbit in the intersection
with $X_2$.
If $O$ itself belongs to $U_2$, evaluation identifies the
simply transitive group of C.36 with $O$ and its universal
cover with $X$. The lifted translations are the unique
based source lifts used here; their transitive generated
group is exactly $N$. The two fractional-power formulas
are consequently identical on $O$ and its closure.
An open orbit not in $U_2$ is disjoint from $X_2$, because
it is an open neighborhood disjoint from $U_2$.
These observations prove agreement on the entire
intersection with $X_2$.

For $h\in H$ with $hO=O'$, lift $h|_O$ to a source-cover
homeomorphism $\widetilde h:X\to X'$.
Covering uniqueness for the two based isotopies gives
\[
              \widetilde h\,\widetilde f_t\,
                 \widetilde h^{-1}=\widetilde f'_t .
                                                        \tag{C.49.16}
\]
Conjugation is consequently a continuous group isomorphism
$N\to N'$. It intertwines their exponential maps and
the unique logarithms in (C.49.8).
It therefore carries (C.49.11) to the corresponding
homotopy on $O'$. Equality extends to the closures.
Changing the lift of $h$ by a deck map does not change
this identity, since deck maps commute with the generated
lifted groups.
The actual boundary-quotient families are transported
at the same time: with basepoint $hx$, their exact
formula is
$\Phi^{O'}_{hz}(lhx)=h\Phi^O_z(lx)$.
Thus the class of open orbits under consideration is
$H$-invariant; no arbitrary selected triple is assumed
to be preserved by the operators.

Every closure and $X_2$ is closed in the plane.
The matching families, and separately their matching
inverse families, paste continuously over a finite union.
Each preserves its member closures and is inverted by
the pasted inverse, so it is a homeomorphism of that union.
These assertions include all original covering and
operator labels. Uniform control for infinitely many
new homothety orbits, and extension to the remaining
one- and two-foliation regions, are not supplied by
this finite-pasting argument.
\end{proof}
